\documentclass[11pt]{article}
\usepackage[T1]{fontenc}
\usepackage[utf8]{inputenc}
\usepackage{amsmath,amssymb,amsthm}
\usepackage{mathtools}
\usepackage{geometry}
\geometry{margin=1.2in}
\usepackage[colorlinks=true,linkcolor=blue,citecolor=blue]{hyperref}
\usepackage{enumitem}

\newtheorem{theorem}{Theorem}[section]
\newtheorem{proposition}[theorem]{Proposition}
\newtheorem{lemma}[theorem]{Lemma}
\newtheorem{corollary}[theorem]{Corollary}
\newtheorem{conjecture}[theorem]{Conjecture}
\theoremstyle{definition}
\newtheorem{definition}[theorem]{Definition}
\newtheorem{example}[theorem]{Example}
\newtheorem{construction}[theorem]{Construction}
\theoremstyle{remark}
\newtheorem{remark}[theorem]{Remark}

\newcommand{\RHom}{\mathrm{RHom}}
\newcommand{\Ext}{\mathrm{Ext}}
\newcommand{\Hom}{\mathrm{Hom}}
\newcommand{\End}{\mathrm{End}}
\newcommand{\Sym}{\mathrm{Sym}}
\newcommand{\Tr}{\mathrm{Tr}}
\newcommand{\HH}{\mathrm{HH}}
\newcommand{\HC}{\mathrm{HC}}
\newcommand{\CE}{\mathrm{CE}}
\newcommand{\Ainf}{A_\infty}
\newcommand{\Linf}{L_\infty}
\newcommand{\Rep}{\mathrm{Rep}}
\newcommand{\Coh}{\mathrm{Coh}}
\newcommand{\Fuk}{\mathrm{Fuk}}
\newcommand{\Perf}{\mathrm{Perf}}
\newcommand{\CoHA}{\mathrm{CoHA}}
\newcommand{\Ran}{\mathrm{Ran}}
\newcommand{\Fact}{\mathrm{Fact}}
\newcommand{\cA}{\mathcal{A}}
\newcommand{\cB}{\mathcal{B}}
\newcommand{\cC}{\mathcal{C}}
\newcommand{\cF}{\mathcal{F}}
\newcommand{\cO}{\mathcal{O}}
\newcommand{\cZ}{\mathcal{Z}}
\newcommand{\frakg}{\mathfrak{g}}
\newcommand{\frakgl}{\mathfrak{gl}}

\title{M-Theory on Calabi--Yau Threefolds, Brane Algebras, \\
  and Koszul Duality \\[8pt]
  \large Physics Notes for Volume III}
\author{Raeez Lorgat}
\date{April 2026}

\begin{document}

\maketitle

\tableofcontents

\bigskip

\noindent
These notes spell out the physical picture behind the quantum vertex chiral
group $G(X)$ of a Calabi--Yau threefold $X$.  The four ingredients ---
M-theory holography, open string field theory from CY categories, closed
strings from cyclic homology, and the packaging of both into a single algebraic
object via Koszul duality --- are developed in turn.  Throughout, we connect to
the rigorous constructions of the monograph: the CY-to-chiral functor $\Phi$
(Theorem CY-A), the bar-cobar machine of Volume~I, and the $E_1/E_2$ chiral
hierarchy of the present Volume~III.


%% ====================================================================
\section{M-theory setup and the two algebras}
\label{sec:mtheory-setup}
%% ====================================================================

\subsection{The geometry}

Fix a compact (or suitably non-compact) Calabi--Yau threefold $X$.  Consider
M-theory on the eleven-dimensional background
\[
  \mathbb{R}_t \times X \times \mathbb{R}^4,
\]
where $\mathbb{R}_t$ is the ``time'' direction along which operator products
will be taken.  Place a stack of $N$ M2-branes supported on
\[
  \mathbb{R}_t \times \{p\} \subset \mathbb{R}_t \times X,
\]
where $p \in X$ is a point.  The M2-branes are thus (1+0)-dimensional objects
(quantum-mechanical from the brane worldvolume perspective) wrapping the time
direction and localized at a point in $X$.

\begin{remark}
The choice of $\mathbb{R}^4$ as the transverse factor is dictated by the
requirement that the M2-brane worldvolume theory be topological along $X$ and
that the bulk gravitational theory reduce to a tractable form after
compactification on $X$.  In the topological string setting, one replaces
$\mathbb{R}^4$ by $\mathbb{C}^{0|2}$ (the odd cotangent directions of the
Calabi--Yau).  This replacement is the AKSZ origin of the BV formalism on the
brane.
\end{remark}

\subsection{The brane algebra $\cA_N$}

The stack of $N$ M2-branes supported on $\mathbb{R}_t \times \{p\}$ carries a
worldvolume theory.  The algebra of local operators of this quantum-mechanical
system is the \emph{brane algebra}
\[
  \cA_N = \text{algebra of operators on the $N$-brane stack at $p \in X$}.
\]
The operator product expansion (OPE) in the $\mathbb{R}_t$ direction equips
$\cA_N$ with the structure of an associative algebra.  More precisely, $\cA_N$
is a differential graded (dg) algebra in the BV formalism, with $Q$-cohomology
the physical operator algebra.

\begin{remark}[Explicit form for $X = \mathbb{C}^3$]
\label{rem:c3-explicit}
When $X = \mathbb{C}^3$ (or more generally a toric CY3), the brane worldvolume
theory is a matrix quantum mechanics.  In the BV formalism, the fields are
\[
  \alpha \in \Omega^*(\mathbb{R}_t)[\varepsilon_1, \varepsilon_2, \varepsilon_3]
    \otimes \frakgl_N[1],
\]
where the $\varepsilon_i$ are odd variables representing the three normal
directions in $X$ (the Ext algebra $\Ext^*_{\cO_X}(\cO_p, \cO_p) =
\bigwedge(\varepsilon_1, \varepsilon_2, \varepsilon_3)$ for a point $p \in
\mathbb{C}^3$).  The action is
\[
  S(\alpha) = \int_{\mathbb{R}_t \times \mathbb{C}^{0|3}}
    \left( \frac{1}{2} \Tr(\alpha \, d\alpha)
    + \frac{1}{3} \Tr(\alpha^3) \right),
\]
where the integral over $\mathbb{C}^{0|3}$ is the Berezin integral over the
odd fiber.  De-BV-fying (extracting the ghost-number-zero fields), one obtains
a gauge field $A \in \Omega^1(\mathbb{R}_t) \otimes \frakgl_N$ and three
$\frakgl_N$-valued scalars $\phi_1, \phi_2, \phi_3$ with action
\[
  S = \int_{\mathbb{R}_t} \Tr\bigl(
    \phi_1 \, d_A \phi_2
    + \phi_3 [\phi_1, \phi_2]
  \bigr).
\]
This is the dimensional reduction of holomorphic Chern--Simons theory on
$\mathbb{C}^3$ to $\mathbb{R}_t$, i.e., a matrix model with cubic potential
$W = \Tr(XYZ - XZY)$.
\end{remark}

\subsection{The gravitational algebra $\cB$}

The \emph{gravitational algebra} $\cB$ is the algebra of local operators of the
topological/holomorphic gravitational theory on the full spacetime
$\mathbb{R}_t \times X$.  Concretely, after compactification on $X$, this is
the BCOV-type theory whose fields are polyvector fields on $X$ (subject to
the $\partial$-closedness condition), and whose OPE in $\mathbb{R}_t$ provides
the associative product.

For $X = \mathbb{C}^3$ (or a neighborhood thereof), the gravitational theory is
the Kodaira--Spencer theory:
\begin{align}
  \text{Fields:} \quad
  &\alpha \in \Omega^*(\mathbb{R}_t) \,\hat{\otimes}\,
    \Omega^{0,*}(X)[1], \notag \\
  &\beta \in \Omega^*(\mathbb{R}_t) \,\hat{\otimes}\,
    \Omega^{0,*}(X)[2], \notag \\[4pt]
  \text{Action:} \quad
  &S = \int_{\mathbb{R}_t \times X}
    \beta \, \bar{\partial} \alpha \, \Omega
    + \beta \, \partial\alpha \wedge \partial\alpha, \label{eq:ks-action}
\end{align}
where $\Omega$ is the holomorphic volume form.  The second term encodes the
Poisson bracket on Hamiltonian vector fields: $\partial\alpha \wedge
\partial\alpha = \{\alpha, \alpha\} \, \Omega$.

The ring of classical local operators is
\[
  \cB^{\mathrm{cl}} = C^*\bigl(
    \mathbb{C}[z_1, z_2, z_3] \ltimes \mathbb{C}[z_1, z_2, z_3][1]
  \bigr),
\]
the Chevalley--Eilenberg cochains of the semi-direct product of the polynomial
ring (viewed as a Lie algebra via the Poisson bracket $\{-,-\}$ from $\Omega$)
with itself shifted, viewed as the adjoint module.

\subsection{Holography as Koszul duality}
\label{sec:holography-koszul}

The holographic duality between the brane theory and the gravitational theory
admits a sharp algebraic formulation:

\begin{conjecture}[Holographic Koszul duality]
\label{conj:holographic-kd}
In the large-$N$ limit, the brane algebra and the gravitational algebra are
related by Koszul duality:
\[
  \lim_{N \to \infty} \cA_N \;=\; \cB^{!},
\]
where $\cB^!$ denotes the Koszul dual of $\cB$.
\end{conjecture}

\begin{remark}[Evidence in the $\mathbb{C}^2$ example]
The original notes verify this for $X = \mathbb{C}^2$ (mixed A-B model on
$\mathbb{R}^2_A \times \mathbb{C}^2_B$).  There, the gauge theory operators
are $\Sym^*(\mathbb{C}[z_1, z_2] \oplus \mathbb{C}[z_1, z_2][1])$ and the
gravity operators are $C^*(\mathbb{C}[z_1, z_2] \ltimes \mathbb{C}[z_1,
z_2][1])$.  The Koszul dual of the latter is $\mathbf{U}(\Hom(\mathbb{C}^2)
\oplus \Hom(\mathbb{C}^2)[1])$, whose commutative (classical, large-$N$) limit
is the former.  Koszul duality thus interchanges:
\begin{center}
\begin{tabular}{rcl}
  Classical free gravity & $\longleftrightarrow$ & Classical gauge theory \\
  Poisson bracket (interactions) & $\longleftrightarrow$ &
    First quantum correction ($\hbar^1$) \\
  Genus-$g$ string amplitude & $\longleftrightarrow$ &
    $\hbar^g$ correction (ribbon graphs)
\end{tabular}
\end{center}
The $A_\infty$-operations in the gravitational Chevalley--Eilenberg complex
(Feynman diagrams with $m_2, m_3, \ldots$) are matched by multi-trace
operators and their OPE on the gauge theory side.
\end{remark}

\begin{remark}[Nature of the limit]
The ``$\lim_{N \to \infty}$'' should be understood as follows.  For each
finite $N$, the brane algebra $\cA_N$ involves $\frakgl_N$-valued fields.
Invariant theory provides the algebra of gauge-invariant operators
\[
  \cA_N^{\mathrm{inv}} = \bigl\{
    \Tr(P(\phi_1, \ldots, \phi_d)) : P \text{ non-commutative polynomial}
  \bigr\} / \text{trace relations at rank } N.
\]
As $N \to \infty$ the trace relations disappear and the generators become free:
$\cA_\infty^{\mathrm{inv}}$ is the free algebra on single-trace generators.
This free algebra is Koszul dual to the Chevalley--Eilenberg cochains of the
gravitational Lie algebra.
\end{remark}


%% ====================================================================
\section{Open strings from CY categories}
\label{sec:open-strings}
%% ====================================================================

\subsection{The CY category as the category of branes}

The data of M-theory on $\mathbb{R}_t \times X$ includes a category of branes:
\begin{itemize}
  \item In the B-model, branes are coherent sheaves on $X$:
    $\cC = D^b(\Coh(X))$.
  \item In the A-model, branes are Lagrangians:
    $\cC = \Fuk(X)$.
  \item In mixed models (e.g., $X = \mathbb{R}^2_A \times \mathbb{C}^2_B$),
    branes are products $L \times Z$ of Lagrangians and coherent sheaves.
\end{itemize}
In every case, $\cC$ is a \emph{Calabi--Yau category of dimension $d = \dim_\mathbb{C} X$}: it carries a cyclic $\Ainf$-structure induced by the holomorphic volume form $\Omega \in H^0(X, K_X)$.

\subsection{Open string fields as $\RHom$}

Fix a brane $\cF \in \cC$ (for instance, the skyscraper sheaf $\cO_p$ for M2-branes at a point $p$, or a Lagrangian $L$ for an A-brane).  The space of \emph{open string fields} is
\[
  \boxed{\;
    \text{Open string fields} = \RHom_\cC(\cF, \cF)[1].
  \;}
\]
The shift $[1]$ is the BV shift: it places the connection (gauge field) in
ghost number zero and the matter fields (normal deformations of the brane) in
appropriate ghost numbers.

\begin{example}[Point brane on $\mathbb{C}^3$]
For $\cF = \cO_p$ with $p \in \mathbb{C}^3$, the self-Ext algebra is
\[
  \Ext^*_{\cO_{\mathbb{C}^3}}(\cO_p, \cO_p)
    = \bigwedge(\varepsilon_1, \varepsilon_2, \varepsilon_3),
\]
an exterior algebra on three generators of degree $1$ (one per normal
direction).  After the BV shift and tensoring with $\frakgl_N$ for $N$ branes,
this gives the field content of Remark~\ref{rem:c3-explicit}.
\end{example}

\begin{example}[Point brane on $\mathbb{C}^2$]
For the mixed A-B model on $\mathbb{R}^2 \times \mathbb{C}^2$, the brane
$\cF = \mathbb{R} \times \{p\}$ has open string fields
\[
  \Omega^*(\mathbb{R})
    \otimes \Ext^*_{\cO_{\mathbb{C}^2}}(\cO_p, \cO_p)
    \otimes \frakgl_N[1]
  = \Omega^*(\mathbb{R})[\varepsilon_1, \varepsilon_2]
    \otimes \frakgl_N[1],
\]
the Floer--Ext tensor product.
\end{example}

\subsection{The open string field theory action from $\Ainf$-structure}
\label{sec:osft-action}

The CY category $\cC$ carries an $\Ainf$-structure: higher composition maps
$\mu_n \colon \Hom(X_0, X_1) \otimes \cdots \otimes \Hom(X_{n-1}, X_n)
\to \Hom(X_0, X_n)[2-n]$.  For the endomorphism algebra $A = \End_\cC(\cF)$
of a single brane, the $\Ainf$-operations $\mu_1, \mu_2, \mu_3, \ldots$
determine the open string field theory action by the Witten--Gaberdiel--Zwiebach
prescription:
\[
  S(\alpha) = \sum_{n \geq 1} \frac{1}{n+1}
    \langle \alpha, \mu_n(\alpha, \ldots, \alpha) \rangle,
\]
where $\langle \cdot, \cdot \rangle \colon A \otimes A \to k[-d]$ is the CY
pairing and $\alpha \in A[1]$ is the string field.

\begin{remark}[Relation to Theorem CY-A]
\label{rem:osft-theorem-cya}
The passage
\[
  \text{cyclic } \Ainf\text{-algebra } (A, \langle \cdot, \cdot \rangle)
  \;\longmapsto\;
  \text{chiral algebra } \Phi(\cC)
\]
of Theorem CY-A is \emph{precisely} the open string field theory construction
in categorical language.  The cyclic $\Ainf$-structure on $A = \End_\cC(\cF)$
is the chain-level datum encoding the full tree-level open string field theory.
The CY-to-chiral functor $\Phi$ promotes this to a factorization algebra on a
curve $C$ by:
\begin{enumerate}[label=(\roman*)]
  \item Extracting the Lie conformal algebra $\mathfrak{L}_\cC$ from the
    Gerstenhaber bracket on $\HH^*(\cC)$;
  \item Taking the factorization envelope $\Fact_C(\mathfrak{L}_\cC)$;
  \item Enhancing to $E_2$ via the $\mathbb{S}^d$-framing (for $d = 3$, the
    $\mathbb{S}^3$-framing gives $E_3$, which restricts to $E_2$ via Dunn
    additivity);
  \item Quantizing using the CY trace.
\end{enumerate}
The upshot: the chiral algebra $\Phi(\cC)$ is the open string field theory
\emph{organized as a factorization algebra}.
\end{remark}

\subsection{$N$ branes and matrix enlargement}

For a stack of $N$ branes, one replaces the endomorphism algebra by its matrix
enlargement:
\[
  A_N = \End_\cC(\cF^{\oplus N}) \simeq \End_\cC(\cF) \otimes \frakgl_N.
\]
The $\Ainf$-structure extends by the Leibniz rule on the $\frakgl_N$ factor.
The brane algebra $\cA_N$ is the $Q$-cohomology of the resulting dg algebra
of gauge-invariant operators:
\[
  \cA_N = H^*\Bigl(
    \bigl(\Sym^*(A_N[1])\bigr)^{\frakgl_N}, \, Q_{\mathrm{BV}}
  \Bigr).
\]
Concretely, generators are single-trace operators $\Tr(P(\phi_1, \ldots,
\phi_d))$ with $P$ a non-commutative polynomial in the matter fields (and
ghosts/anti-fields at other ghost numbers).  As $N \to \infty$, the trace
relations vanish and $\cA_\infty$ is freely generated by single traces.


%% ====================================================================
\section{Closed strings from cyclic homology}
\label{sec:closed-strings}
%% ====================================================================

\subsection{The theorem of Loday--Quillen--Tsygan}

The passage from open strings to closed strings is controlled by a classical
theorem in non-commutative geometry:

\begin{theorem}[Loday--Quillen--Tsygan]
\label{thm:lqt}
For a unital associative algebra $A$, the Chevalley--Eilenberg homology of the
Lie algebra of infinite matrices over $A$ is an exterior algebra generated by
the cyclic homology of $A$:
\[
  H_*^{\CE}(\frakgl_\infty(A)) \;\simeq\; \bigwedge \HC_*(A).
\]
\end{theorem}

\subsection{Physical interpretation}
\label{sec:lqt-physics}

The physical content of Loday--Quillen--Tsygan is as follows:

\begin{enumerate}[label=(\arabic*)]
  \item \textbf{The open string algebra} is $A = \End_\cC(\cF)$, the
    endomorphism algebra of a brane.

  \item \textbf{Matrix enlargement and the large-$N$ limit.}  The Lie algebra
    $\frakgl_N(A) = \frakgl_N \otimes A$ is the open string algebra for $N$
    branes.  Taking $N \to \infty$ yields $\frakgl_\infty(A)$.

  \item \textbf{Chevalley--Eilenberg homology = gauge-invariant operators.}
    The CE chain complex $C_*^{\CE}(\frakgl_\infty(A))$ computes the homology
    of gauge-invariant observables in the large-$N$ matrix model.  This is the
    space of \emph{closed string states}: the gauge-invariant sector of the
    open string theory IS the closed string theory.

  \item \textbf{Single-trace generators = single-string states.}  The exterior
    algebra generators $\HC_*(A)$ correspond to single-trace operators
    $\Tr(\phi_{i_1} \cdots \phi_{i_k})$.  These are the \emph{single-string
    states}: each cyclic homology class is one closed string mode.

  \item \textbf{Multi-trace operators = multi-string states.}  The exterior
    power $\bigwedge^k \HC_*(A)$ corresponds to $k$-trace operators, i.e.,
    $k$-string states.  The exterior algebra structure encodes the statistics
    (closed strings are bosonic/fermionic depending on degree).

  \item \textbf{OPE = string scattering.}  The OPE of single-trace operators
    on the gauge theory side matches the tree-level closed string scattering
    amplitudes.  Multi-trace OPE (with trace-splitting and trace-joining)
    matches higher-genus string amplitudes.
\end{enumerate}

\subsection{The closed string algebra as chiral derived center}
\label{sec:derived-center}

In the language of the monograph:

\begin{proposition}[Closed strings = derived center]
\label{prop:closed-derived-center}
The closed string algebra produced by Loday--Quillen--Tsygan is the
\emph{chiral derived center} $Z^{\mathrm{der}}_{\mathrm{ch}}(\Phi(\cC))$ of
the chiral algebra associated to the CY category:
\[
  \text{Closed string algebra}
  = H_*^{\CE}(\frakgl_\infty(A))
  \simeq \bigwedge \HC_*(A)
  = Z^{\mathrm{der}}_{\mathrm{ch}}(\Phi(\cC)).
\]
\end{proposition}

The identifications proceed as follows:
\begin{itemize}
  \item $\HC_*(A)$ is the cyclic homology of the open string algebra ---
    equivalently, the Hochschild homology of $\cC$ with its $S^1$-action.
  \item $\HC_*(A)$ is simultaneously the space of deformations of the closed
    string background (via the Kontsevich--Soibelman map $\HC_*(A) \to$
    deformations of the $\Ainf$-structure).
  \item The chiral derived center $Z^{\mathrm{der}}_{\mathrm{ch}}(\Phi(\cC))$
    is the $E_2$-center of the chiral algebra --- exactly the algebra of bulk
    (gravitational) local operators.
\end{itemize}

\begin{remark}
The identification of the derived center with the gravitational algebra is the
categorical incarnation of the Drinfeld center construction:
\[
  \cZ(\Rep^{E_1}(\Phi(\cC)))
  \simeq \Rep^{E_2}(Z^{\mathrm{der}}_{\mathrm{ch}}(\Phi(\cC))).
\]
The monoidal ($E_1$) representation category of the boundary (brane) theory
determines a braided monoidal ($E_2$) category via the center construction.
This is the bulk-boundary correspondence at the categorical level.
\end{remark}


%% ====================================================================
\section{The gravitational algebra from polyvector fields}
\label{sec:grav-algebra}
%% ====================================================================

\subsection{BCOV theory on $X$}
\label{sec:bcov}

The gravitational (closed string field) theory on $\mathbb{R}_t \times X$ is
the Kodaira--Spencer / BCOV theory.  Following the analysis in the earlier
notes, and extending from $\mathbb{C}^2$ to a general CY threefold $X$:

The space of closed string states (before interaction) is
\[
  \ker \partial \subset \mathbf{PV}^{*,*}(X) \simeq
  \bigoplus_{p,q} \Omega^{0,q}(X, \bigwedge^p TX),
\]
where $\partial \colon \mathbf{PV}^{p,q} \to \mathbf{PV}^{p-1,q}$ is the
divergence operator (contraction with $\Omega$).  The components:
\begin{itemize}
  \item $\mathbf{PV}^{0,*}(X) \subset \ker \partial$: functions (automatic).
  \item $\mathbf{PV}^{1,*}(X) \cap \ker \partial$: divergence-free vector
    fields, i.e., Hamiltonian vector fields (via $\Omega$, these are identified
    with closed holomorphic $(d-1)$-forms).
  \item $\mathbf{PV}^{2,*}(X) \cap \ker \partial$: $\partial$-closed
    bi-vector fields.  For $X = \mathbb{C}^3$, these give deformations of the
    holomorphic symplectic structure.
  \item $\mathbf{PV}^{3,*}(X) \cap \ker \partial$: the ``top'' polyvector
    fields, corresponding to deformations of $\Omega$ itself.
\end{itemize}

\subsection{The gravitational algebra as Chevalley--Eilenberg cochains}
\label{sec:grav-ce}

After the reduction to physical operators (passing to $Q$-cohomology), the
ring of classical local operators of the gravitational theory at a point
$x \in X$ is
\[
  \cB^{\mathrm{cl}}
  = C^*\bigl(\frakgl_\Omega(X) \ltimes \frakgl_\Omega(X)[1]\bigr),
\]
where $\frakgl_\Omega(X) = \{f \in \cO(X) : \{f, -\} \text{ preserves }
\Omega\}$ is the Lie algebra of Hamiltonian functions with Poisson bracket
$\{-,-\}$ determined by $\Omega$.  The semi-direct product structure reflects
the BF-type nature of the theory: $\alpha$-fields (Hamiltonian vector fields)
act on $\beta$-fields (adjoint-valued functions) via the Poisson bracket, while
$\beta$ commutes with $\beta$.

For computations at a formal neighborhood of $p \in X \simeq \mathbb{C}^3$:
\begin{align*}
  \text{Ghost \#1 operators:} &\quad
    \cO_i^{(1)} = \partial_{z_1}^{a} \partial_{z_2}^{b} \partial_{z_3}^{c}
    \alpha \quad (a, b, c \geq 0), \\
  \text{Ghost \#2 operators:} &\quad
    \cO_i^{(2)} = \partial_{z_1}^{a} \partial_{z_2}^{b} \partial_{z_3}^{c}
    \beta \quad (a, b, c \geq 0).
\end{align*}
These transform \emph{contravariantly} to the gauge theory operators: the
gravitational operators carry polynomial derivatives (jet space), while the
gauge theory operators carry polynomial values (formal functions).  This
contravariance is the fingerprint of Koszul duality.


%% ====================================================================
\section{The Koszul duality dictionary}
\label{sec:koszul-dictionary}
%% ====================================================================

\subsection{Classical Koszul duality}

At the classical level, the Koszul duality between gauge theory and gravity
takes the form:

\begin{center}
\renewcommand{\arraystretch}{1.4}
\begin{tabular}{|c|c|c|}
\hline
& \textbf{Gauge theory (brane)} & \textbf{Gravity (bulk)} \\
\hline
\textbf{Algebra} & $\cA_\infty = \Sym^*\bigl(V \oplus V[1]\bigr)$ &
  $\cB = C^*(V \ltimes V[1])$ \\
\hline
\textbf{Type} & Commutative (symmetric) & Lie (CE cochains) \\
\hline
\textbf{Generators} & $\Tr(\phi^I)$ (polynomial values) &
  $\partial^I \alpha, \, \partial^I \beta$ (jets) \\
\hline
\textbf{Product} & Commutative product & Chevalley--Eilenberg cup product \\
\hline
\end{tabular}
\end{center}

\noindent
where $V = \mathbb{C}[z_1, \ldots, z_d]$ (polynomial functions on the
transverse CY directions).  The Koszul dual of $C^*(V \ltimes V[1])$ is
$\mathbf{U}(V \oplus V[1])$ --- the universal enveloping algebra --- whose
commutative limit (setting $\hbar \to 0$, equivalently $N \to \infty$ with
classical scaling) is $\Sym^*(V \oplus V[1])$.

\subsection{Quantum corrections: order by order in $\hbar$}
\label{sec:quantum-corrections}

Moving beyond the classical limit, Koszul duality becomes a precise matching
between quantum corrections:

\begin{enumerate}[label=\textbf{Order $\hbar^{\arabic*}$:}]
  \item[{$\hbar^0$:}]
    Classical free gravity $\longleftrightarrow$ Classical gauge theory
    (commutative limit).

  \item[{$\hbar^1$:}]
    The Poisson bracket $\{z_1^k z_2^l, z_1^m z_2^n\}$ on the gravity side
    matches the commutator $[\Tr(\phi_1^k \phi_2^l), \Tr(\phi_1^m \phi_2^n)]$
    on the gauge theory side.  Diagrammatically, this is a single propagator
    connecting two single-trace sub-diagrams.

  \item[{$\hbar^2$:}]
    By the even quantization symmetry ($\hbar \mapsto -\hbar$ with reversal of
    product), there are no corrections at order $\hbar^2$.

  \item[{$\hbar^3$:}]
    The non-trivial $A_\infty$-operations appear.  On the gauge theory side,
    ribbon graphs with $k$ boundaries contribute $k$-trace operators.  A cubic
    ribbon graph (three boundaries) gives an $m_3$ operation; a quadratic one
    gives a corrected $m_2$.  On the gravity side, these match Feynman diagrams
    computing the $A_\infty$-structure on $C^*(\frakg_\Omega)$.
\end{enumerate}

\begin{remark}[Genus expansion]
The $\hbar$ expansion is simultaneously a genus expansion.  At order $\hbar^g$,
the diagrams on the gauge theory side are ribbon graphs of genus $\leq g$; on
the gravity side, they are genus-$g$ string worldsheets.  The precise statement
is:
\begin{itemize}
  \item $\hbar^0$: tree level (genus 0, planar diagrams)
  \item $\hbar^1$: one-loop (torus diagrams, first non-planar correction)
  \item $\hbar^g$: genus-$g$ worldsheets $\longleftrightarrow$ ribbon graphs
    with $g$ handles
\end{itemize}
In the Volume~I language, this is the shadow obstruction tower: genus-$g$ contributions are
captured by the arity-$(g+1)$ component of the MC element $\Theta_A$.
\end{remark}


%% ====================================================================
\section{The quantum vertex chiral group $G(X)$}
\label{sec:qvcg}
%% ====================================================================

\subsection{Packaging the two algebras}

The quantum vertex chiral group $G(X)$ associated to a CY3 $X$ is a single
algebraic object that packages both the brane algebra and the gravitational
algebra, with Koszul duality as the internal organizing principle:

\begin{definition}[Quantum vertex chiral group --- physics summary]
\label{def:gx-physics}
The quantum vertex chiral group $G(X)$ is the $E_2$-chiral algebra
$\Phi(\cC)$, where $\cC$ is the CY category of $X$ (either $D^b(\Coh(X))$ or
$\Fuk(X)$), equipped with the following structure:
\begin{enumerate}[label=(\roman*)]
  \item \textbf{Positive half (brane algebra).}  The $E_1$-sector of $G(X)$ is
    the CoHA: $G(X)^+ = \CoHA(\cC) \simeq \cA_\infty$.  For toric $X$, this
    is $Y^+(\widehat{\frakg}_{Q_X})$ (the positive half of the affine super
    Yangian).

  \item \textbf{Derived center (gravitational algebra).}  The chiral derived
    center $Z^{\mathrm{der}}_{\mathrm{ch}}(G(X)) = \cB$ is the gravitational
    algebra.  It is generated as an exterior algebra by $\HC_*(\cC)$
    (Loday--Quillen--Tsygan).

  \item \textbf{Koszul duality = holography.}  The positive half and the
    derived center are related by Koszul duality:
    \[
      G(X)^+ \simeq \bigl(Z^{\mathrm{der}}_{\mathrm{ch}}(G(X))\bigr)^!.
    \]
    This is the algebraic incarnation of holographic duality.

  \item \textbf{$E_2$-structure (braiding = $R$-matrix).}  The $E_2$-chiral
    enhancement provides the braiding on $\Rep^{E_2}(G(X))$.  The universal
    $R$-matrix $R(z) = \mathrm{Res}^{\mathrm{coll}}_{0,2}(\Theta_{G(X)})$
    satisfies the quantum Yang--Baxter equation and encodes the quantum group
    structure.
\end{enumerate}
\end{definition}

\subsection{The triangular decomposition}

As a super Lie algebra (or BKM superalgebra in the K3 $\times$ E setting),
$G(X)$ admits a triangular decomposition:
\[
  G(X) = \mathfrak{n}_- \oplus \mathfrak{h} \oplus \mathfrak{n}_+,
\]
where:
\begin{itemize}
  \item $\mathfrak{n}_+ = $ positive half $=$ CoHA $=$ brane algebra
    $\cA_\infty$.
  \item $\mathfrak{h} = $ Cartan subalgebra $=$ lattice $\Lambda(X) \otimes
    \mathbb{R}$ extracted from $K$-theory of $X$.
  \item $\mathfrak{n}_- = $ negative half $=$ co-brane algebra (Koszul dual
    generators).
\end{itemize}
The root space decomposition
\[
  \mathfrak{n}_+ = \bigoplus_{\alpha \in \Delta_+}
    \mathfrak{n}_\alpha, \qquad
  \dim \mathfrak{n}_\alpha = \mathrm{mult}(\alpha)
\]
is controlled by the BPS spectrum: $\mathrm{mult}(\alpha) = \mathrm{DT}_\alpha(X)$
(Donaldson--Thomas invariant of class $\alpha$).  The real roots correspond to
rigid BPS states; the imaginary roots to families.

\subsection{The automorphic correction and the shadow obstruction tower}
\label{sec:automorphic-shadow}

The passage from the ``tree-level'' root datum (real roots + leading imaginary
roots) to the full BKM superalgebra is the \emph{automorphic correction}.  In
the language of Volume~I:

\begin{center}
\renewcommand{\arraystretch}{1.4}
\begin{tabular}{|c|c|}
\hline
\textbf{BKM / physics} & \textbf{Volume I bar-cobar} \\
\hline
Tree-level $r$-matrix & Collision $r$-matrix
  $r(z) = \mathrm{Res}^{\mathrm{coll}}_{0,2}(\Theta_A)$ \\
\hline
Imaginary root corrections & Higher-arity shadows
  ($C$, $Q$, \ldots) \\
\hline
Full automorphic correction & Full MC element $\Theta_A$ \\
\hline
Denominator identity $\Phi(z)$ & Bar-complex Euler product \\
\hline
Automorphic form (e.g., $\Delta_5$) & Genus-$g$ obstruction $\mathrm{obs}_g$ \\
\hline
\end{tabular}
\end{center}

The key identification:
\begin{center}
\emph{Automorphic correction = Shadow obstruction tower.}
\end{center}
The shadow Postnikov obstruction tower $\Theta_A^{\leq r}$ of Volume~I, projected at
successive arities, recovers the successive imaginary root corrections to the
BKM superalgebra.

\subsection{The two main examples}
\label{sec:two-examples}

\subsubsection{Example 1: $X = \mathbb{C}^3$ (affine Yangian)}

\begin{itemize}
  \item Quiver: Jordan quiver (one vertex, one loop with cubic potential
    $W = \Tr(XYZ - XZY)$).
  \item $\cA_\infty = \CoHA(\mathbb{C}^3) \simeq Y^+(\widehat{\frakgl}_1)$
    (Schiffmann--Vasserot).
  \item $G(\mathbb{C}^3) = Y(\widehat{\frakgl}_1)$, the full affine Yangian
    of $\frakgl_1$.
  \item This is isomorphic to $\mathcal{W}_{1+\infty}$ at the self-dual level
    --- already a central object in Volume~I, where MC4 is solved by weight
    stabilization.
  \item $\cB = $ Kodaira--Spencer theory on $\mathbb{C}^3$ = the chiral
    derived center.
  \item The Yangian $R$-matrix is the DK-0 shadow.
\end{itemize}

\subsubsection{Example 2: $X = (\text{K3} \times E) / (\mathbb{Z}/N\mathbb{Z})$
  (BKM superalgebras)}

\begin{itemize}
  \item Lattice: $\Lambda^{3,2}$ of signature $(3,2)$ from $H^*(X)$.
  \item Real roots: three simple roots $\delta_1, \delta_2, \delta_3$ with Gram
    matrix
    $\bigl(\begin{smallmatrix} 2 & -2 & -2 \\ -2 & 2 & -2 \\ -2 & -2 & 2
    \end{smallmatrix}\bigr)$.
  \item Root multiplicities: $\mathrm{mult}(\alpha) = f(nm, l)$, where
    $f(n, l)$ are Fourier coefficients of the K3 elliptic genus
    $\varphi_{0,1}$.
  \item $G(X) = \frakg_{\Delta_5}$, the BKM superalgebra whose denominator
    identity is (up to normalization) the Igusa cusp form $\Delta_5$, a Siegel
    modular form of weight 5 for $\mathrm{Sp}_4(\mathbb{Z})$.
  \item The $E_2$-structure is the $\mathrm{Sp}_4(\mathbb{Z})$ action on the
    genus-2 Siegel upper half-space $\mathbb{H}_2$.
  \item The shadow obstruction tower captures the automorphic correction: arity 2 is the
    Weyl vector contribution, higher arities add imaginary roots.
\end{itemize}


%% ====================================================================
\section{The full picture: open/closed/holography}
\label{sec:full-picture}
%% ====================================================================

We can now assemble the complete dictionary:

\[
\boxed{
\begin{array}{ccccc}
  \textbf{Open strings} & \xrightarrow{\text{LQT}} &
  \textbf{Closed strings} & \xleftrightarrow{\text{KD}} &
  \textbf{Branes} \\[6pt]
  \RHom(\cF, \cF)[1] & &
  \bigwedge \HC_*(\End(\cF)) & &
  \cA_N \\[6pt]
  \Ainf\text{-structure} & &
  Z^{\mathrm{der}}_{\mathrm{ch}} & &
  \text{CoHA} \\[6pt]
  E_1\text{-sector} & &
  E_2\text{-bulk} & &
  E_1\text{-boundary}
\end{array}
}
\]

Reading left to right:

\begin{enumerate}[label=\textbf{(\arabic*)}]
  \item \textbf{Open strings.}  The CY category $\cC$ provides the space of
    open string fields $\RHom(\cF, \cF)[1]$ with action from the $\Ainf$-structure.  This is Theorem CY-A: the cyclic $\Ainf$-structure is promoted to a chiral algebra $\Phi(\cC)$.

  \item \textbf{Closed strings.}  Loday--Quillen--Tsygan extracts the closed
    string algebra as $\bigwedge \HC_*(\End(\cF))$ --- the exterior algebra
    on cyclic homology.  This is the chiral derived center
    $Z^{\mathrm{der}}_{\mathrm{ch}}(\Phi(\cC))$ and equals the gravitational
    algebra $\cB$.

  \item \textbf{Holography.}  The brane algebra $\cA_\infty$ (single-trace
    operators, the CoHA, the positive half of $G(X)$) is Koszul dual to $\cB$.
    The Drinfeld center construction
    $\cZ(\Rep^{E_1}(\cA)) \simeq \Rep^{E_2}(\cB)$
    provides the categorical bulk-boundary correspondence.

  \item \textbf{The quantum vertex chiral group.}  $G(X)$ packages all three:
    $G(X)^+ = \cA_\infty$ (brane algebra / CoHA),
    $Z^{\mathrm{der}}_{\mathrm{ch}}(G(X)) = \cB$ (gravitational algebra),
    and Koszul duality between them is holography.  The $E_2$-structure
    (braiding / $R$-matrix) encodes the quantum group structure, and the full
    MC element $\Theta_{G(X)}$ (the automorphic correction) is the shadow obstruction tower
    capturing all genera.
\end{enumerate}

\subsection{Summary of identifications}

\begin{center}
\renewcommand{\arraystretch}{1.5}
\begin{tabular}{|l|l|l|}
\hline
\textbf{Physics} & \textbf{Algebra} & \textbf{Vol I/III} \\
\hline
M2-brane operators & $\cA_N$ (gauge-invariant operators) & CoHA, $G(X)^+$ \\
\hline
Gravitational operators & $\cB$ (CE cochains) &
  $Z^{\mathrm{der}}_{\mathrm{ch}}(G(X))$ \\
\hline
Open string fields & $\RHom(\cF, \cF)[1]$ & Cyclic $\Ainf$ \\
\hline
Open string action & $\sum \frac{1}{n+1} \langle \alpha, \mu_n(\alpha^n)
  \rangle$ & CY trace + $\Ainf$ \\
\hline
Closed string states & $\HC_*(\End(\cF))$ & Cyclic homology \\
\hline
Multi-string Fock space & $\bigwedge \HC_*(\End(\cF))$ & LQT theorem \\
\hline
Holographic duality & $\cA_\infty = \cB^!$ & Koszul duality \\
\hline
Bulk-boundary & $\cZ(\Rep^{E_1}(\cA)) \simeq \Rep^{E_2}(\cB)$ &
  Drinfeld center \\
\hline
$R$-matrix & QYB solution & $\mathrm{Res}^{\mathrm{coll}}_{0,2}(\Theta)$ \\
\hline
BPS spectrum & DT invariants & Root multiplicities \\
\hline
String scattering & OPE of traces & $\Ainf$-operations via KD \\
\hline
Genus-$g$ amplitude & Ribbon graph, $\hbar^g$ & Shadow obstruction tower arity $g+1$ \\
\hline
\end{tabular}
\end{center}


%% ====================================================================
\section{Closing remarks and open directions}
\label{sec:open-directions}
%% ====================================================================

\begin{enumerate}[label=(\arabic*)]
  \item \textbf{Wall-crossing.}  The BPS spectrum of $X$ is subject to
    wall-crossing (Kontsevich--Soibelman).  In the $G(X)$ language,
    wall-crossing should correspond to gauge equivalence of the MC element
    $\Theta_{G(X)}$: different stability conditions give gauge-equivalent
    $\Ainf$-structures, hence isomorphic quantum vertex chiral groups.  This
    is the content of the $\Ainf$-wall-crossing formula.

  \item \textbf{Beyond CY3.}  For CY $d$-folds with $d \neq 3$, the story
    changes:
    \begin{itemize}
      \item $d = 2$ (K3 surfaces): the $\mathbb{S}^2$-framing gives an
        $E_2$-algebra directly.  No M-theory lift needed; the story is purely
        about 2d string theory.
      \item $d = 4$ (CY 4-folds): the brane is an M5-brane (or the
        topological twist thereof), and the brane algebra is 3-dimensional.
        The relevant algebraic structure is an $E_3$-algebra (via the
        $\mathbb{S}^4$-framing and Dunn).
    \end{itemize}

  \item \textbf{Non-toric CY3.}  For non-toric CY3 (e.g., the quintic), the
    CoHA construction is more subtle: one lacks a quiver presentation, and the
    root datum must be extracted from the full derived category or Fukaya
    category.  The conjecture is that $G(X)$ still exists (as a quantum vertex
    chiral group) but may not admit a BKM or Yangian presentation.

  \item \textbf{Closed string field theory.}  The gravitational algebra $\cB$
    is the \emph{classical} closed string field theory.  The full quantum
    closed string field theory is controlled by the genus expansion, i.e., by
    the full shadow obstruction tower $\Theta_{G(X)}$.  The quantum master equation
    $d\Theta + \frac{1}{2}[\Theta, \Theta] = 0$ is the BV master equation of
    the closed string field theory.

  \item \textbf{Multiple brane types.}  With multiple branes $\cF_1, \ldots,
    \cF_k$, the open string algebra is enriched to the full $\Ainf$-category
    $\cC$ (not just an endomorphism algebra).  The brane algebra becomes a
    factorization algebra on the configuration space of brane types.  For $k$
    brane types with $E_2$ directions, one expects $E_2$-Koszul duality in the
    sense of the braided factorization theory of Part~IV.
\end{enumerate}

\end{document}
